\documentclass[11pt]{article}

\usepackage[margin=1in]{geometry}
\usepackage[T1]{fontenc}
\usepackage[utf8]{inputenc}
\usepackage{amsmath}
\usepackage{amssymb}
\usepackage{array}
\usepackage{booktabs}
\usepackage{float}
\usepackage{tabularx}
\usepackage[hidelinks]{hyperref}

\title{XTSF: Synthetic Forecasting Contributions with XPC\\
Theory, Estimands, Implementation, and Experiments}
\author{Gaspard Berthelier}
\date{\today}

\newcommand{\code}[1]{\texttt{#1}}
\newcommand{\E}{\mathbb{E}}
\newcommand{\R}{\mathbb{R}}
\newcommand{\MSE}{\operatorname{MSE}}
\newcolumntype{Y}{>{\raggedright\arraybackslash}X}

\begin{document}
\maketitle

\begin{abstract}
XTSF studies whether Monte Carlo Shapley explanations recover known
contributions in a controlled forecasting problem.  A synthetic hourly load is
generated as the sum of explicit weather, calendar, trend, and irrelevant
components, so the structural contribution of every group is known.  The
experiments separate two sources of error: Monte Carlo error relative to the
chosen Shapley estimand, and estimand mismatch relative to the structural
decomposition.  They compare interventional and nearest-neighbor conditional
Shapley values, vary the number of sampled coalitions and masking samples, and
reproduce three aggregation games from the original XPC work.  The accompanying
\code{xpc} package provides the reusable data, model, masking, grouping,
estimation, and post-processing abstractions used by the notebook.
\end{abstract}

\setcounter{tocdepth}{1}
\tableofcontents
\clearpage

\section{Research Question and Scope}

The main question is not merely whether a Shapley estimator converges.  It is
whether the quantity to which it converges corresponds to the contribution of
interest.  For an observation \(x\), a forecasting model \(f\), and a known
structural decomposition \(C^{\mathrm{struct}}(x)\), the project asks:

\begin{enumerate}
  \item How quickly does Monte Carlo Shapley converge as the number of
        coalitions and masking samples increases?
  \item How does an interventional coalition value differ from a conditional
        coalition value when features are dependent?
  \item How close is each exact Shapley estimand to the known structural
        contribution?
  \item How do alternative aggregation games trade computation for a changed
        estimand?
\end{enumerate}

The word \emph{contribution} therefore refers to three distinct objects:

\begin{itemize}
  \item the structural contribution specified by the synthetic data-generating
        process;
  \item the exact Shapley contribution for a chosen empirical coalition-value
        function;
  \item its finite-budget Monte Carlo estimate.
\end{itemize}

Keeping these objects separate is essential.  A low error relative to an exact
conditional Shapley value only establishes numerical convergence to that
conditional estimand.  It does not establish agreement with the structural or
causal decomposition.

\section{Theoretical Framework}

\subsection{Shapley Values}

Let \(N=\{1,\ldots,p\}\) be a set of players and let \(v_x(S)\) be the value of
coalition \(S\subseteq N\) for the observation \(x\).  The Shapley value of
player \(j\) is

\begin{equation}
\phi_j(x;v)
=
\sum_{S\subseteq N\setminus\{j\}}
\frac{1}{p\binom{p-1}{|S|}}
\left[v_x(S\cup\{j\})-v_x(S)\right].
\label{eq:shapley}
\end{equation}

The weights in Equation~\eqref{eq:shapley} are equivalent to first drawing the
coalition size uniformly from \(0,\ldots,p-1\), then drawing a subset of that
size uniformly from the remaining players.  For a single coherent game, exact
Shapley values satisfy efficiency:

\begin{equation}
\sum_{j=1}^{p}\phi_j(x;v)=v_x(N)-v_x(\varnothing).
\label{eq:efficiency}
\end{equation}

In the core package, players may be raw features or explicit disjoint feature
groups.  Changing the groups changes the cooperative game; groups are not
merely display labels.

\subsection{Baseline, Interventional, and Conditional Values}

For a deterministic reference vector \(b\), baseline masking defines

\begin{equation}
v_x^{\mathrm{base}}(S)=f(x_S,b_{\bar S}).
\end{equation}

For a random feature vector \(X\), the interventional value is

\begin{equation}
v_x^{\mathrm{int}}(S)
=\E\left[f(x_S,X_{\bar S})\right].
\label{eq:interventional}
\end{equation}

It preserves the dependence among jointly missing features, but breaks their
dependence with the features fixed to \(x_S\).  Consequently, it may evaluate
combinations that are uncommon in the observed data.

The conditional value is

\begin{equation}
v_x^{\mathrm{cond}}(S)
=\E\left[f(x_S,X_{\bar S})\mid X_S=x_S\right].
\label{eq:conditional}
\end{equation}

Conditional masking stays closer to the observed support because missing
features are drawn from rows compatible with the present coalition.  This can
be more realistic for prediction explanations, but it is still associational.
Observed dependence may arise from confounding, common causes, time, or feature
construction.  Conditional Shapley is therefore not automatically a causal
attribution.

\subsection{Exact Empirical Estimands}

Let \(B=\{z^{(1)},\ldots,z^{(R)}\}\) be the finite background data.  The exact
empirical interventional value used in the notebook is

\begin{equation}
\widehat v_x^{\mathrm{int}}(S)
=\frac{1}{R}\sum_{r=1}^{R}f(x_S,z^{(r)}_{\bar S}).
\label{eq:exact-interventional}
\end{equation}

For conditional values, the observed columns \(S\) are standardized by their
background standard deviations.  Let \(\mathcal N_K(x_S)\) be the \(K=64\)
nearest background rows in that standardized space.  The notebook defines

\begin{equation}
\widehat v_x^{\mathrm{cond}}(S)
=\frac{1}{K}\sum_{z\in\mathcal N_K(x_S)}f(x_S,z_{\bar S}),
\label{eq:exact-conditional}
\end{equation}

with all background rows used when \(S=\varnothing\).  An exact empirical
Shapley reference is obtained by evaluating these deterministic coalition
values and enumerating every subset in Equation~\eqref{eq:shapley}.

Here, \emph{exact} means exact for the finite empirical or nearest-neighbor
game.  It removes coalition-sampling and mask-sampling noise.  It is not an
exact population conditional expectation and it is not a causal ground truth.

\section{Synthetic Structural Ground Truth}

\subsection{Data-Generating Process}

The notebook creates 120 days of hourly observations, for \(R=2880\) rows.
Temperature combines a slow seasonal component, an hourly component, and
Gaussian noise.  This induces dependence between weather variables, time of
day, weekday status, and trend.  The six model features and structural groups
are summarized below.

\begin{table}[H]
\centering
\small
\begin{tabularx}{\textwidth}{@{}l l Y r@{}}
\toprule
Group & Feature & Definition or role & Coefficient \\
\midrule
Weather & heating degree & \(\max(18-T,0)\) & 1.20 \\
Weather & cooling degree & \(\max(T-23,0)\) & 1.80 \\
Calendar & daily peak & morning and evening Gaussian basis functions & 14.0 \\
Calendar & weekday work & indicator for Monday through Friday & 6.0 \\
Trend & trend & linear ramp from 0 to 1 & 10.0 \\
Irrelevant & irrelevant & independent standard Gaussian feature & 0.0 \\
\bottomrule
\end{tabularx}
\caption{Synthetic forecasting features and structural coefficients.}
\label{tab:features}
\end{table}

The load and forecasting model are identical and contain no observation noise:

\begin{equation}
f(x)=60
+1.2x_{\mathrm{heat}}
+1.8x_{\mathrm{cool}}
+14x_{\mathrm{peak}}
+6x_{\mathrm{weekday}}
+10x_{\mathrm{trend}}.
\label{eq:model}
\end{equation}

The true group contributions are the corresponding additive terms:

\begin{align}
C_{\mathrm{weather}}(x)
  &=1.2x_{\mathrm{heat}}+1.8x_{\mathrm{cool}},\\
C_{\mathrm{calendar}}(x)
  &=14x_{\mathrm{peak}}+6x_{\mathrm{weekday}},\\
C_{\mathrm{trend}}(x)&=10x_{\mathrm{trend}},\\
C_{\mathrm{irrelevant}}(x)&=0.
\end{align}

Thus the structural decomposition is known for every row.  With zero-baseline
masking, the grouped Shapley values equal these terms exactly because the model
is additive and every basis feature is absent at zero.  This is used as an
oracle sanity check of tensor shapes and grouping behavior.

\subsection{Putting Signed Shapley Values on the Structural Scale}

Empirical Shapley values explain deviations from the empty-coalition value.
The notebook compares them with positive structural contributions by restoring
the matching background mean:

\begin{equation}
\widehat C_g(x)=\widehat\phi_g(x)+\overline C_g,
\qquad
\overline C_g=\frac{1}{R}\sum_{r=1}^{R}C_g(z^{(r)}).
\label{eq:mean-allocation}
\end{equation}

For the additive model, exact interventional Shapley satisfies
\(\phi_g^{\mathrm{int}}(x)=C_g(x)-\overline C_g\), so
Equation~\eqref{eq:mean-allocation} recovers structural truth.  Conditional
Shapley may reallocate signal across dependent groups, so the same mean
allocation does not generally recover each structural component.

Two mean squared errors are reported over evaluation rows \(i\) and groups
\(g\):

\begin{align}
\MSE_{\mathrm{struct}}
  &=\frac{1}{nG}\sum_{i,g}
    \left(\widehat C_{i,g}-C^{\mathrm{struct}}_{i,g}\right)^2,\\
\MSE_{\mathrm{estimand}}
  &=\frac{1}{nG}\sum_{i,g}
    \left(\widehat C_{i,g}-C^{\mathrm{exact}}_{i,g}\right)^2.
\end{align}

The second quantity measures Monte Carlo error.  The gap between the exact
estimand and structural truth is an estimand mismatch that more computation
cannot remove.

\section{Monte Carlo Implementation}

\subsection{Estimator}

For every explained unit and player \(g\), the estimator draws
\(K=\code{n\_coalitions}\) coalitions from the Shapley subset distribution.
For each coalition it estimates

\begin{equation}
\widehat v_x(S)=\frac{1}{M}\sum_{m=1}^{M}
f\!\left(\widetilde x_S^{(m)}\right),
\qquad M=\code{n\_mask\_samples},
\end{equation}

then averages the marginal differences

\begin{equation}
\widehat\phi_g(x)=\frac{1}{K}\sum_{k=1}^{K}
\left[\widehat v_x(S_k\cup\{g\})-\widehat v_x(S_k)\right].
\end{equation}

Increasing \code{n\_coalitions} reduces subset-sampling noise.  Increasing
\code{n\_mask\_samples} reduces noise in the empirical coalition-value
expectation.  The implementation samples the before and after coalition values
independently.

\subsection{Package Abstractions}

The installable package is named \code{xpc}.  Its main responsibilities are:

\begin{table}[H]
\centering
\small
\begin{tabularx}{\textwidth}{@{}l Y@{}}
\toprule
Abstraction & Implemented responsibility \\
\midrule
\code{DataSpec} & Tabular \((n,f)\) data and normalized \((n,H)\) outputs. \\
\code{TimeSeriesTensorSpec} & Temporal \((n,d,f)\) data; masks one time row while the model receives the full sequence, returning \((n,d,H,G)\) explanations. \\
\code{ModelAdapter} & NumPy boundary for callables, pandas, sklearn-style, PyTorch, external-script, R-script, and registered custom models. \\
\code{Masker} & Deterministic baseline, unconditional empirical, or conditional empirical replacement. \\
\code{Conditioner} & Nearest-neighbor, grid-cell, or user-defined non-generative conditional selection. \\
\code{FeatureGroups} & Disjoint named players, remaining-feature policies, and always-present features. \\
\code{ShapleyExplainer} & Monte Carlo subset estimator returning signed values and optional heightened parts. \\
\code{Explanation} & Values, base values, predictions, efficiency residuals, heightening, and comparison metrics. \\
\bottomrule
\end{tabularx}
\caption{Current XPC implementation boundaries.}
\label{tab:abstractions}
\end{table}

For temporal inputs, background rows used by empirical maskers are drawn from
the flattened collection of time locations.  Conditional matching operates on
the features of the currently explained row; it is not a temporal generative
model.

Always-present features are excluded from the player set and restored in every
coalition.  They receive no separate attribution.  They affect the base value
and, in models with interactions, may also change the marginal contributions
assigned to other players.

\subsection{Explanation and Efficiency Diagnostics}

An \code{Explanation} stores signed values \(\Phi\), base values \(b\), and
model predictions \(\widehat y\).  Its efficiency residual is
\(r=\widehat y-(b+\sum_g\widehat\phi_g)\).

For a coherent exact Shapley game this residual is zero by
Equation~\eqref{eq:efficiency}.  In the package it diagnoses finite Monte Carlo
error, independently estimated coalition values, and model-output shape
handling.  Changing from interventional to conditional masking changes the
estimand, but does not by itself invalidate Shapley efficiency.

\subsection{Heightening}

Heightening is optional positive post-processing and is enabled by default.
For every output and group, the implementation computes an offset over all
batch axes,

\begin{equation}
o_{h,g}=\alpha\max\!\left(-\min_u\Phi_{u,h,g},0\right),
\qquad \alpha\geq 0,
\end{equation}

then forms non-negative values, percentages, and target allocations:

\begin{align}
\Phi^+_{u,h,g}&=\max(\Phi_{u,h,g}+o_{h,g},0),\\
p_{u,h,g}&=\frac{\Phi^+_{u,h,g}}{\sum_j\Phi^+_{u,h,j}},\\
\mathrm{part}_{u,h,g}&=p_{u,h,g}T_{u,h}.
\end{align}

The target defaults to the non-negative model prediction.  If all positive
values are zero, a uniform allocation is used.  Heightening leaves the signed
values untouched.  All convergence and aggregation MSE experiments use signed
values with \code{heighten=False}; heightened values are visualization outputs,
not the estimand being benchmarked.

\section{Convergence Experiment}

The convergence benchmark uses the four structural groups in the direct
grouped game.  The full 2880-row panel is the background.  Six evaluation rows
are selected at evenly spaced indexes, and every setting is repeated five
times with deterministic but distinct random seeds.

Two one-dimensional sweeps are performed:

\begin{itemize}
  \item \code{n\_coalitions} in \(\{1,2,4,8,16,32\}\), with
        \code{n\_mask\_samples}=16;
  \item \code{n\_mask\_samples} in \(\{1,2,4,8,16,32\}\), with
        \code{n\_coalitions}=16.
\end{itemize}

For interventional Shapley, the additive structural contribution is also the
exact empirical estimand after the mean allocation in
Equation~\eqref{eq:mean-allocation}.  For conditional Shapley, the notebook
enumerates all coalitions and uses Equation~\eqref{eq:exact-conditional} to
construct a separate exact empirical reference.

As computation increases, the interventional structural MSE tends toward zero.
Conditional MSE relative to its own estimand also decreases, but its structural
MSE plateaus near the exact conditional estimand gap of \(7.31\).  For example,
at 32 coalitions and 16 mask samples, the conditional MSE is \(0.172\) relative
to its estimand and \(7.305\) relative to structural truth.  This is the central
distinction of the experiment: convergence does not remove estimand bias.

\section{Aggregation Games}

The archived XPC work considered a climate group \(C\) and a non-climate group
\(D=N\setminus C\).  The notebook reproduces three definitions.

\subsection{Default Post-hoc Aggregation}

First compute a Shapley value for every raw feature, then sum values belonging
to the group:

\begin{equation}
\phi_C^{\mathrm{post}}(x)=\sum_{j\in C}\phi_j(x;v).
\end{equation}

This preserves the efficiency of the original raw-feature game after all
groups are summed, but requires explaining every raw feature.

\subsection{Coalitional Shapley}

For each requested group \(C\), merge all features in \(C\) into one player
while every feature outside \(C\) remains an individual player.  With
\(q_C=p-|C|+1\),

\begin{equation}
\phi_C^{\mathrm{coal}}(x)
=\sum_{S\subseteq N\setminus C}
\frac{1}{q_C\binom{q_C-1}{|S|}}
\left[v_x(S\cup C)-v_x(S)\right].
\label{eq:coalitional}
\end{equation}

Climate and non-climate are computed in two different reduced games because
their complements contain different raw players.  Each value is a valid
Shapley value in its own game, but the pair is not generally jointly efficient.
This corrects the stronger efficiency claim made in the historical report.

\subsection{Simplified Two-player Shapley}

Treat \(C\) and \(D\) as the only two players in one common game:

\begin{align}
\phi_C^{\mathrm{two}}(x)
  &=\tfrac12\left([v_x(C)-v_x(\varnothing)]
                 +[v_x(N)-v_x(D)]\right),\\
\phi_D^{\mathrm{two}}(x)
  &=\tfrac12\left([v_x(D)-v_x(\varnothing)]
                 +[v_x(N)-v_x(C)]\right).
\end{align}

Only the four coalition values
\(v(\varnothing),v(C),v(D),v(N)\) are needed.
The pair is efficient because both values belong to the same two-player game.

The package's direct grouped-player API generalizes this idea to any number of
explicit groups through \code{ShapleyExplainer}.  The exact historical post-hoc,
coalitional, and two-player comparison is implemented locally in the notebook,
not as an \code{aggregation} parameter in the core package.

\section{Aggregation Experiment and Results}

The fixed-budget comparison uses 16 sampled coalitions, 16 masking samples per
coalition value, six evaluation rows, and five repetitions.  The simplified
mode does not sample coalitions because the two-player formula is exhaustive;
its four coalition values are still estimated from 16 masked rows each.

\begin{table}[H]
\centering
\small
\begin{tabular}{@{}llrrrr@{}}
\toprule
Masking & Mode & Exact structural & MC structural & MC estimand & Rows/point \\
\midrule
Interventional & Post-hoc & 0.000 & 0.512 & 0.512 & 3072 \\
Interventional & Coalitional & 0.000 & 0.171 & 0.171 & 1024 \\
Interventional & Two-player & 0.000 & 1.594 & 1.594 & 64 \\
Conditional & Post-hoc & 11.398 & 12.478 & 1.362 & 3072 \\
Conditional & Coalitional & 10.943 & 10.718 & 0.573 & 1024 \\
Conditional & Two-player & 10.203 & 11.572 & 1.267 & 64 \\
\bottomrule
\end{tabular}
\caption{Aggregation-mode mean squared errors.  ``Exact structural'' is the
exact empirical estimand's MSE against structural truth.  Monte Carlo columns
are means over five repetitions.}
\label{tab:aggregation-results}
\end{table}

The masked-row count follows directly from the implemented loops.  Post-hoc
uses \(6\times16\times2\times16=3072\) model rows per point.  Coalitional uses
\(2\times16\times2\times16=1024\).  The two-player mode uses
\(4\times16=64\), a factor of 48 below post-hoc aggregation.  Measured timings
show the same ordering but are machine-dependent.

All interventional exact aggregation estimands coincide with structural truth
because the model is additive.  Conditional exact estimands differ because
conditioning reallocates signal among dependent features and because the three
aggregation definitions are different games.  The smallest conditional exact
structural MSE in this sample is not evidence that the corresponding mode is
generally more causal.  It is a property of this data, model, grouping, and
nearest-neighbor estimand.

The exact conditional Coalitional pair has an efficiency RMSE of \(0.842\),
while post-hoc and the common two-player game are efficient up to numerical
precision.  This confirms that independently computed Coalitional group values
should not be interpreted as a single efficient decomposition.

\section{Outputs and Reproducibility}

The single Colab-ready notebook is \code{src/xtsf.ipynb}.  It caches the
synthetic panel and writes generated artifacts under \code{outputs/}:

\begin{itemize}
  \item \path{outputs/synthetic_forecasting_panel.csv};
  \item \path{outputs/shapley_convergence_runs.csv};
  \item \path{outputs/shapley_convergence_summary.csv};
  \item \path{outputs/shapley_convergence.png};
  \item \path{outputs/aggregation_exact_estimands.csv};
  \item \path{outputs/aggregation_mode_runs.csv};
  \item \path{outputs/aggregation_mode_summary.csv};
  \item \path{outputs/aggregation_modes.png}.
\end{itemize}

The package implementation is under \code{src/xpc/}; lightweight tests are
under \code{src/tests/}.  The public estimator deliberately excludes known
contributions, SMACH-specific processing, GAM-specific paths, Kernel SHAP,
Permutation SHAP, and generative conditional models.  External or known
contributions enter through \code{Explanation.from\_contributions} and
\code{Explanation.compare}, not through the estimator.

\section{Interpretation and Limitations}

\begin{itemize}
  \item The exact structural decomposition is available because the synthetic
        model is additive and matches the data-generating process.  More complex
        models with interactions will not make the three interventional
        aggregation estimands coincide.
  \item Conditional nearest-neighbor masking avoids many unsupported feature
        combinations, but remains a finite-sample associational approximation.
        It does not identify interventions or causal effects.
  \item Exact empirical references remove Monte Carlo noise but retain finite
        background, neighbor-selection, model, and grouping choices.
  \item The benchmarks use six evaluation points and five repetitions.  They
        demonstrate estimand and computation effects, but are not a definitive
        large-sample performance study.
  \item Heightened positive parts are useful forecast allocations, but they are
        post-processed visual quantities.  Scientific error comparisons use the
        signed Shapley values plus the documented background-mean allocation.
\end{itemize}

\end{document}
